\documentclass[11pt]{article}
\usepackage[margin=1in]{geometry}
\usepackage{amsmath,amssymb,booktabs}
\begin{document}

\begin{center}
{\large\bfseries Study reference: forward-LLP setup and provided data}
\end{center}

\author{Tony Menzo}

\noindent
This reference collects the fixed inputs for the forward sensitivity study
defined in the task. The physics model is stated in the task prompt, and the
decays are to be derived from it. Everything here is the part of the setup that
is \emph{not} derived: the toy experimental configuration and the provided
production data. All values are synthetic and are to be used exactly as given.

\section*{Experimental configuration}

\begin{center}
\begin{tabular}{@{}lll@{}}
\toprule
Quantity & Value & Role \\
\midrule
Beam & $pp$, $\sqrt{s} = 13.6$~TeV & set in \texttt{pythia\_cards/} \\
$\sigma_{\mathrm{inel}}$ & $80.0$~mb & inelastic $pp$ cross section \\
$N_{\mathrm{int}}$ & $2.7\times10^{18}$ & number of primary interactions \\
Sensitivity rule & $N_{\mathrm{sig}} = 3$, zero background & defines the reach \\
Geometry & \texttt{geometry/geometry.yaml} & decay volume, detector plane, acceptance \\
\bottomrule
\end{tabular}
\end{center}

\noindent
Two searches are defined, a two-track dimuon search and a collinear diphoton
search, with acceptances specified in the geometry file. Each mass is evaluated
under the search whose decay channel is open there. How the generated flux is
normalized and combined into a yield at $N_{\mathrm{int}}$ interactions is part
of the analysis and is not prescribed here.

\section*{Forward flux generation}

The forward parent flux is generated with the two Pythia cards in
\texttt{pythia\_cards/}:
\begin{itemize}
\item \texttt{forward\_flux.cmnd} --- soft inelastic QCD, which produces the
  light mesons and open charm ($K$, $D$, $D_s$, $\rho$, $\omega$, $\phi(1020)$).
\item \texttt{forward\_onia.cmnd} --- prompt charmonium ($J/\psi$, $\psi(2S)$)
  through NRQCD.
\end{itemize}

\section*{Provided production}

Production of $\phi$ is provided as a data product; the production matrix
elements are not to be re-derived. For each open channel, per parent meson:
\begin{itemize}
\item \texttt{production/kappa.csv} --- the $g^2$-stripped branching
  $\kappa = \mathrm{Br}(\text{parent}\to\cdots\,\phi)$ at $g=1$, per
  $(\text{parent}, m_\phi)$. A parent contributes $\phi$ with per-collision
  weight $w_{\text{parent}}\,\kappa\,g^2$, where $w_{\text{parent}}$ is its
  weight in the forward flux.
\item \texttt{spectra/spectrum\_$\langle$parent$\rangle$\_m$\langle m_\phi\rangle$.csv}
  --- the $\phi$ energy spectrum in the parent rest frame, tabulated as
  $(x, \mathrm{pdf})$ with $x = 2E_\phi^\star/m_{\text{parent}}$ (shape only).
  To place one $\phi$: draw $x$ from this distribution, take the $\phi$
  direction isotropic in the parent rest frame, and boost to the lab with the
  parent momentum.
\end{itemize}

\section*{Production channels and accessible mass range}

The scalar is radiated from a parent meson only below the endpoint of the
$(\phi + X)$ phase space, where $X$ collects the other final-state particles,
\[
m_\phi \;<\; m_{\text{parent}} - \sum_{i\,\in\,X} m_i .
\]
For the two production topologies, with the neutrino taken massless, this is
\[
M \to \mu\,\nu\,\phi:\quad m_\phi < m_M - m_\mu ,
\qquad
V \to \mu^+\mu^-\phi:\quad m_\phi < m_V - 2 m_\mu .
\]
The resulting per-parent endpoints are:

\begin{center}
\begin{tabular}{@{}lccc@{}}
\toprule
Parent & Topology & $m_{\text{parent}}$ [GeV] & $m_\phi$ endpoint [GeV] \\
\midrule
$K$          & $M\to\mu\nu\phi$ & $0.494$ & $0.388$ \\
$D$          & $M\to\mu\nu\phi$ & $1.870$ & $1.764$ \\
$D_s$        & $M\to\mu\nu\phi$ & $1.968$ & $1.863$ \\
$\rho$       & $V\to\mu\mu\phi$ & $0.775$ & $0.564$ \\
$\omega$     & $V\to\mu\mu\phi$ & $0.783$ & $0.571$ \\
$\phi(1020)$ & $V\to\mu\mu\phi$ & $1.019$ & $0.808$ \\
$J/\psi$     & $V\to\mu\mu\phi$ & $3.097$ & $2.886$ \\
$\psi(2S)$   & $V\to\mu\mu\phi$ & $3.686$ & $3.475$ \\
\bottomrule
\end{tabular}
\end{center}

\noindent
The accessible range extends up to the heaviest parent's endpoint. The study is
evaluated on the pinned mass grid [GeV]
\[
0.06,\ 0.12,\ 0.18,\ 0.205,\ 0.23,\ 0.28,\ 0.34,\ 0.42,\ 0.70,\ 1.20,\ 1.90,\ 2.30,
\]
and the rows present in \texttt{kappa.csv} are exactly the channels open at each
mass.

\section*{Probe points}

Report the following scalar probes:

\begin{center}
\begin{tabular}{@{}lll@{}}
\toprule
Field & Probe point & Quantity \\
\midrule
\texttt{width\_probe\_gev} & $m_\phi=0.30,\ g=4.2\times10^{-4}$ & $\Gamma(\phi\to\mu\mu)$ [GeV] \\
\texttt{width\_gg\_probe\_gev} & $m_\phi=0.12,\ g=1.0\times10^{-3}$ & $\Gamma(\phi\to\gamma\gamma)$ [GeV] \\
\texttt{ctau\_probe\_m[0]} & $m_\phi=0.30,\ g=3.7\times10^{-6}$ & $c\tau$ [m] (total width) \\
\texttt{ctau\_probe\_m[1]} & $m_\phi=0.30,\ g=1.11\times10^{-5}$ & $c\tau$ [m] (total width) \\
\bottomrule
\end{tabular}
\end{center}

\end{document}
